\section{Analysis: When Does Folding Help?}
\label{sec:analysis}

The benchmark ladder separates three recurring regimes. On BCIC2020-3, a
native 64-row input and full visual fine-tuning outperform every displayed EEG
foundation model, suggesting that architecture and supervised adaptation are
already strong. On TUAB and SEED-V, the completed visual runs trail the
strongest EEG reference. FACED supplies a third regime: ConvNeXt-Tiny exceeds
the EEG references while ViT-Small is substantially lower, showing sensitivity
to the backbone and training recipe. The analyses below connect these outcomes to fold factor and channel
neighborhoods rather than treating the 12 tasks as an undifferentiated ranking.

\subsection{Fold factor controls the geometry seen by the CNN}

The visual aspect ratio before padding is

\begin{equation}
  \frac{W}{H}=\frac{T}{CP^2}.
\end{equation}

Increasing $P$ changes this ratio quadratically, but a square image is not the
objective.  Large $P$ also makes each horizontal step span more original
samples and increases the distance across phase-wrap boundaries.  Consistent with this trade-off, Eq.~\ref{eq:p-geometry-rule} provides a
geometry-based default from $(C,T)$ and the EfficientNet-B0 stride. It is not
expected to be universally optimal because task-relevant EEG structure can be
short-range or long-range in time and local or distributed across channels.
Table~\ref{tab:p-rule} gives the completed unified choices.

\begin{table*}[t]
\centering
\caption{Architecture-aligned fold factors used in the unified reproduction.}
\label{tab:p-rule}
\small
\begin{tabularx}{\textwidth}{@{}lrrrcX@{}}
\toprule
Dataset & $C$ & $T$ & $P_{\mathrm{geo}}$ & Folded $H\times W$ & Status \\
\midrule
CHB-MIT & 16 & 2000 & 4 & $64\times500$ & completed \\
TUAB & 16 & 2000 & 4 & $64\times500$ & completed \\
TUEV & 16 & 1000 & 4 & $64\times250$ & completed \\
ISRUC & 6 & 6000 & 12 & $72\times500$ & completed \\
FACED & 32 & 2000 & 2 & $64\times1000$ & completed \\
SEED-V & 62 & 200 & 1 & $62\times200$ & completed \\
PhysioNet-MI & 64 & 800 & 1 & $64\times800$ & completed \\
SHU-MI & 32 & 800 & 2 & $64\times400$ & completed \\
BCIC2020-3 & 64 & 600 & 1 & $64\times600$ & completed \\
Mumtaz2016 & 19 & 1000 & 4 & $76\times250$ & completed \\
MentalArithmetic & 20 & 1000 & 4 & $80\times250$ & completed \\
HMC & 4 & 6000 & 16 & $64\times375$ & completed \\
\bottomrule
\end{tabularx}
\tablenote{The rule uses only input shape and backbone stride. All listed
factors match the canonical YAML recipes.}
\end{table*}

The rule leaves native 64-channel inputs such as BCIC2020-3 and PhysioNet-MI
unfolded, assigns $P=4$ to the 16-channel recordings, and assigns $P=12$ to
the six-channel ISRUC input. On SEED-V, the canonical EfficientNet-B0
$P=1$ recipe reaches $\kappa=.2056\pm.0078$. This benchmark result does not
establish whether another fold factor would improve performance. Sampling
rate, event duration, channel topology, and task-specific temporal support
may affect the preferred geometry. The MentalArithmetic sweep in
Appendix~\ref{app:p-ablation} provides a concrete example: the best primary
metric occurs at $P=4$ for EfficientNet-B0 but at $P=2$ for both DINOv3
backbones. This supports treating $P$ as a task- and backbone-dependent
hyperparameter when additional validation compute is available.

\paragraph{Two complementary mechanisms.}
Folding is not merely a height-expansion device. The bijection preserves all
waveform values while rearranging them into phase--time neighborhoods, so the
model can detect temporal patterns through standard 2D kernels. Its benefit
therefore combines occupancy with pattern accessibility: it preserves more
signal-bearing height and exposes local relations that were distant or
anisotropic in the original channel--time matrix. Which mechanism dominates is
task dependent and should be tested with folding and permutation controls.

\subsection{Where artificial neighborhoods fail}

Folding preserves channel order but does not encode channel names or 3D
locations.  Permuting electrodes changes the model even if the underlying set
of waveforms is identical.  This position sensitivity can be beneficial when a
fixed montage is used consistently, but it limits transfer across acquisition
setups.  In contrast, \reve{} assigns explicit physical coordinates and is
designed for montage changes.

The next model should therefore keep the cheap folding operator while adding a
minimal topology signal: a channel-ID embedding, coordinate-valued planes, or a
small channel mixer before folding.  A decisive experiment is a channel-order
stress test: train with one order, evaluate under known permutations, and
compare row-order encoding with physical-coordinate encoding.
